\begin{document}

Verallgemeinerte Eigenvektoren berechnen: \\
Setze $v_2 = B \cdot v_3$: 
\[
\begin{pmatrix}
2 & -3 & 1 & -1 \\
3 & -6 & 0 & -3 \\
-1 & 3 & 1 & 2 \\
-3 & 6 & 0 & 3
\end{pmatrix} \cdot \begin{pmatrix}
a\\b\\c\\d
\end{pmatrix} = \begin{pmatrix}
-1 \\ -1 \\ 0 \\ 1
\end{pmatrix} \quad \Rightarrow \quad v_3 = \begin{pmatrix}
0 \\ -\frac{1}{3} \\ 0 \\ 1
\end{pmatrix}
\]
%\txc{Warum setzt man hier $v_3 = B \cdot v_4$? Ich verstehe, dass das irgendetwas mit der Verkettung zu tun hat, aber mehr auch nicht.}
Setze $v_3 = B \cdot v_4$: 
\[
\begin{pmatrix}
2 & -3 & 1 & -1 \\
3 & -6 & 0 & -3 \\
-1 & 3 & 1 & 2 \\
-3 & 6 & 0 & 3
\end{pmatrix} \cdot \begin{pmatrix}
a\\b\\c\\d
\end{pmatrix} = \begin{pmatrix}
0 \\ -\frac{1}{3} \\ 0 \\ 1
\end{pmatrix} \quad \Rightarrow \quad v_3 = \begin{pmatrix}
0 \\ -\frac{1}{3} \\ 0 \\ 1
\end{pmatrix}
\]
Und keine Ahnung, wie es weiter geht oder wo ich $v_4$ her bekomme. \\
Aber ich weiß zumindestens, dass die Jordan Blöcke 2, 1 und 1 groß sein sollten. Also sollte $J$ so aussehen: 
\[
J = \begin{pmatrix}

\end{pmatrix}
\]

\end{document}